% ECONOMICS_PROBLEM: {"id": "K42-506", "title": "Youth recruitment into organized crime and expectations", "jel_code": "K42", "source_id": "OP-0628"}
% FIRST_POSTED: 2026-10-09
% ATTEMPT_STATUS: partial
% ENTRY_KIND: research_attempt
\documentclass[11pt,a4paper]{article}
\usepackage[T1]{fontenc}
\usepackage[utf8]{inputenc}
\usepackage[margin=27mm]{geometry}
\usepackage{amsmath,amssymb,amsthm,mathtools}
\usepackage[hidelinks]{hyperref}
\setlength{\parindent}{0pt}
\setlength{\parskip}{5pt}
\newtheorem{theorem}{Theorem}[section]
\newtheorem{proposition}[theorem]{Proposition}
\newtheorem{lemma}[theorem]{Lemma}
\newcommand{\R}{\mathbb R}
\newcommand{\E}{\mathbb E}
\DeclareMathOperator{\Var}{Var}
\DeclareMathOperator{\Cov}{Cov}
\begin{document}
\section*{Youth recruitment prevention: interventions, beliefs, and network effects}
\section{What a randomized policy can and cannot identify}
The target cohort is fixed before a prevention intervention and followed to a common age. This attempt establishes identification of separately randomized information and opportunity packages, supplies an exact decomposition of their interaction, and derives a bounded linear network benchmark. It also constructs an observational equivalence showing why total effects and post-policy beliefs alone do not identify a belief-mediated effect. The actual effects on gang entry and lawful earnings remain unresolved without data.

Let $I,O\in\{0,1\}$ denote an information package and an opportunity package. Randomize whole noninterfering recruitment clusters to the four policy cells, with positive known assignment probabilities. Outcomes may depend arbitrarily on assignments within their own cluster, but here the cluster package is the complete intervention. For a baseline youth define gang entry $G(I,O)$ by a fixed age, and cumulative lawful earnings $Y(I,O)$ with a declared zero-after-death and zero-during-nonemployment convention. The latter is a composite earnings outcome, not earnings conditional on survival.

Under consistency, random assignment independent of the four potential outcomes, and complete follow-up, cell means identify $m_{io}=E[G(i,o)]$ and the analogous earnings means. If cluster sizes differ, the analysis must use baseline youth weights to target a youth-average effect; unweighted cluster means target a different population. Independent clusters, not individual youth observations, are the independent sampling units for variance and concentration calculations.

\section{A complete intervention decomposition}
The joint package effect is $m_{11}-m_{00}$. Its factorial interaction is
\[
 \iota=m_{11}-m_{10}-m_{01}+m_{00}.
\]
Two order-specific decompositions are valid:
\[
 m_{11}-m_{00}=(m_{10}-m_{00})+(m_{11}-m_{10})
 =(m_{01}-m_{00})+(m_{11}-m_{01}).
\]
A symmetric allocation assigns
\[
 \phi_I=\tfrac12[(m_{10}-m_{00})+(m_{11}-m_{01})],\qquad
 \phi_O=\tfrac12[(m_{01}-m_{00})+(m_{11}-m_{10})]. \tag{1}
\]
Expanding terms proves $\phi_I+\phi_O=m_{11}-m_{00}$. Each component averages the intervention's incremental effect across the two possible orders. This is an explicit allocation convention for package effects, not identification of a psychological mediator.

For example, let $(m_{00},m_{10},m_{01},m_{11})=(0.4,0.3,0.2,0.05)$. The total entry reduction is $0.35$, interaction is $-0.05$, and $(\phi_I,\phi_O)=(-0.125,-0.225)$. The information intervention's effect depends on whether opportunities are also changed. Calling the first component ``the expectations mechanism'' would require proving that the information package has no direct effect on social support, monitoring, or actual opportunities.

\section{A nonidentification witness for natural mediation}
Even a perfectly randomized binary policy $Z$ and a perfectly measured binary belief $M$ need not identify a natural indirect effect. Consider two models with $M(z)=z$ for every youth. Define potential entry outcomes conditional on imposed policy and mediator as follows:
\[
 \text{Model A: }G(z,m)=1-m,\qquad
 \text{Model B: }G(z,m)=1-z.
\]
Both models produce exactly $M=Z$ and $G=1-Z$ under random assignment, so the joint observed law of $(Z,M,G)$ agrees. The total effect is $-1$ in each. But the natural indirect contrast at baseline policy is
\[
 E[G(0,M(1))-G(0,M(0))]=-1\text{ in A},\quad0\text{ in B}. \tag{2}
\]
The calculation evaluates the displayed functions. Model A makes the effect run through beliefs; Model B makes it direct. Their unobserved cross-policy mediator interventions differ. Recording expectations more accurately does not distinguish the models. The deterministic mediator also lacks within-policy overlap, highlighting an additional obstacle; no positivity condition should be silently presumed from policy randomization.

A mediator experiment that directly randomizes a well-defined information state while holding the opportunity intervention fixed would identify a controlled effect under its own exclusion and consistency assumptions. It generally targets a different quantity from (2). Sequential ignorability, mediator overlap and absence of exposure-induced mediator-outcome confounding are possible alternative identifying restrictions, but they require justification rather than follow from the total-policy trial.

\section{A solvable network propagation benchmark}
For a finite network, take a nonnegative adjacency-weight matrix $A$, baseline entry propensity vector $a$, direct prevention shift $bz$ with $b\ge0$, and peer coefficient $\gamma\ge0$. The benchmark expected propensities satisfy
\[
 q(z)=a-bz+\gamma A q(z). \tag{3}
\]
Assume $\gamma\|A\|_\infty<1$ and that the resulting vector lies in $[0,1]^n$ for every intervention considered. The probability restriction is separate; a linear system need not automatically define valid probabilities.
\begin{proposition}
Equation (3) has a unique solution and the prevention response is
\[
 q(z)-q(0)=-b(I-\gamma A)^{-1}z
 =-b\sum_{k=0}^{\infty}\gamma^kA^kz. \tag{4}
\]
For $z\ge0$, every coordinate weakly decreases. If $A\mathbf1=\mathbf1$, full coverage changes every coordinate by $-b/(1-\gamma)$.
\end{proposition}
\begin{proof}
The Neumann series converges in the matrix infinity norm, multiplies $I-\gamma A$ to the identity, and has nonnegative entries. This establishes existence, uniqueness and signs in (4). If $A\mathbf1=\mathbf1$, every power preserves $\mathbf1$, and the geometric series gives the final expression.
\end{proof}
This explicitly displays indirect effects absent from an isolated individual comparison. If $A$ and $\gamma$ are known but $b$ is interval-valued, (4) propagates that interval exactly for nonnegative assignments. If the residual error in (3) has norm at most $\epsilon$, the solution error is at most $\epsilon/(1-\gamma\|A\|_\infty)$. Near the stability boundary, small misspecification can dominate the direct effect. Recruitment displacement, rival-group substitution and endogenous network rewiring are excluded; any of them can invalidate the monotone sign. The benchmark concerns prevention effects, not a prescription for recruitment.

\section{Mortality and observation}
By the fixed horizon, report mutually exclusive states such as entry before death or horizon, death without prior recorded entry, and survival without entry. Death after entry belongs to the first state; survival is also worth reporting separately. Treating death as ordinary missing entry data can make a harmful mortality increase look like successful prevention. The earnings convention should remain the same across regimes. Conditioning on survivors changes the target and introduces principal-stratum selection.

For bounded $G\in[0,1]$ with observation rate $r_z$ and observed mean $g_z$, the full-cohort entry mean lies in $[r_zg_z,r_zg_z+1-r_z]$. These bounds follow by assigning all missing outcomes zero or one. They do not justify substituting zero for missing records. Administrative validation of entry and mortality can narrow observation uncertainty, while misclassification requires a separate sensitivity analysis.

\section{Remaining gap}
The factorial design identifies packages and their interaction; the counterexample blocks an unsupported mediation interpretation; the network equations give a conditional propagation theorem. A full solution needs real cluster assignments, implementation fidelity, independent measurement of entry and survival, and evidence about changes in expectations and opportunities. Different policies can produce the same short-run reported beliefs while changing long-run lawful earnings differently. No factual youth recruitment rate or policy effect is asserted here.
\begin{thebibliography}{9}
\bibitem{crime} S. Tob\'on, M. M. Sviatschi and N. Cabra-Ruiz (2025), \emph{Organised Crime}, VoxDevLit 17(1), September. \url{https://voxdev.org/voxdevlit/organised-crime}. Source motivation for expectations, recruitment and networks; the experimental and mathematical benchmarks above are self-contained.
\end{thebibliography}

\end{document}
